\documentclass[11pt]{article}
\usepackage[utf8]{inputenc}
\usepackage{amsmath,amssymb}
\usepackage{graphicx}
\usepackage{booktabs}
\usepackage{hyperref}
\usepackage[margin=1in]{geometry}
\usepackage{natbib}
\bibliographystyle{unsrtnat}

\title{A Hippocampus-Accumbens Code Guides Goal-Directed Appetitive Behavior}
\author{}
\date{}

\begin{document}
\maketitle

\begin{abstract}
The hippocampal-accumbens circuit is implicated in transforming spatial representations into motivation-driven actions, yet the specific information carried by hippocampal neurons projecting to the nucleus accumbens (dHPC$\to$NAc) during goal-directed behavior remains unknown. Here we use dual-colour two-photon calcium imaging in head-fixed mice performing a spatial reward learning task to simultaneously record from large populations of dHPC$\to$NAc and non-projecting (dHPC$^{-}$) neurons. We find that dHPC$\to$NAc neurons contain a larger proportion of place cells with enhanced spatial information, greater trial-to-trial reliability, and stronger in-field activity. Their place fields over-represent the reward zone and are more strongly modulated by local tactile cues. Speed-inhibited (deceleration) neurons and lick-excited neurons are both overrepresented in the dHPC$\to$NAc population. Optogenetic stimulation of dHPC terminals in the NAc induces mouth movements and deceleration, confirming functional relevance. A generalized linear model reveals enhanced conjunctive coding of position, velocity, and appetitive licking in dHPC$\to$NAc neurons, providing a multiplexed signal that improves reward zone identification. These results demonstrate that the hippocampal output to the accumbens carries an enriched, behaviourally relevant representation that integrates spatial, locomotor, and consummatory information for goal-directed navigation.
\end{abstract}

\section{Introduction}

The hippocampus is renowned for its role in spatial navigation and memory, encoding the animal's position through place cells whose firing is confined to specific locations in the environment \citep{okeefe1971hippocampus,moser2008place}. However, the hippocampus does not act in isolation: its spatial representations must be transmitted to downstream targets that translate them into appropriate behavioural outputs. The nucleus accumbens (NAc), situated at the interface of limbic and motor circuits, is a prime candidate for this transformation \citep{mogenson1980hippocampal,ito2008dissociation}.

Previous work has established that lesioning or silencing hippocampal projections to the NAc diminishes conditioned place preference and impairs spatial reward learning \citep{britt2012synaptic,lefort2019hippocampal}. However, these loss-of-function approaches cannot reveal what information the projection carries. The hippocampus encodes not only position but also velocity, reward expectation, and other task-relevant variables \citep{mcnaughton2006path,gauthier2018dedicated}. Which of these signals are selectively routed to the NAc, and whether NAc-projecting neurons carry enriched or specialized representations compared to the general hippocampal population, has not been determined.

Technical limitations have constrained prior investigations. Electrophysiological recordings identify projection neurons through antidromic activation, but this approach yields small samples of identified cells. Recent advances in dual-colour two-photon imaging now enable simultaneous recording from hundreds of neurons while identifying projection targets through retrograde viral labelling \citep{chen2013ultrasensitive,dana2019high}.

Here we combine retrograde viral tracing with dual-colour two-photon imaging to characterize the coding properties of dHPC$\to$NAc neurons during a head-fixed spatial reward learning task, complemented by optogenetic activation and computational modelling.

\section{Results}

\subsection{Mice learn a spatial reward task}

Food-restricted mice ($n = 18$) were trained on a head-fixed spatial reward learning task using a 360-cm self-propelled treadmill belt with six textured zones and one hidden 30-cm reward zone (Figure~1A). Mice learned across five training days to lick a spout in the reward zone to receive liquid reward. Repeated-measures ANOVA confirmed significant increases in rewarded laps per session ($F_{4} = 6.344$, $p < 0.01$, Greenhouse-Geisser corrected), anticipatory licking in the 30-cm zone preceding the reward ($F_{4} = 3.803$, $p < 0.05$), and reward zone licking ($F_{4} = 4.276$, $p < 0.05$, Greenhouse-Geisser corrected; Figure~1B--D).

\subsection{Dual-colour imaging identifies dHPC$\to$NAc projection neurons}

In six Thy1-GCaMP6s mice, retrograde AAVrg-pgk-Cre was injected into the medial NAc and Cre-dependent DIO-mCherry into dHPC (CA1/subiculum border), enabling dual-colour two-photon imaging: dynamic GCaMP6s calcium signals from all neurons plus static mCherry labelling of NAc-projecting neurons (Supplementary Figure~2A). Spectral separation between the green (GCaMP6s) and red (mCherry) channels was confirmed.

\subsection{dHPC$\to$NAc neurons carry enhanced spatial information}

dHPC$\to$NAc neurons contained a significantly larger proportion of place cells compared to the dHPC$^{-}$ population (Figure~2). Among identified place cells, dHPC$\to$NAc neurons exhibited higher spatial information (bits/event), greater trial-to-trial reliability, and stronger in-place-field calcium activity. These results indicate that the hippocampal output to the accumbens is not a random sample of the hippocampal population but is enriched for neurons with robust, reliable spatial representations.

\subsection{Reward zone overrepresentation and cue modulation}

Place field analysis revealed that dHPC$\to$NAc neurons showed enhanced overrepresentation of the reward zone compared to dHPC$^{-}$ neurons (Figure~3A). Additionally, dHPC$\to$NAc place fields were more strongly modulated by local belt cues: place field edges clustered near the boundaries between textured zones, indicating stronger alignment with tactile landmarks (Figure~3B). This cue-modulated reward zone overrepresentation provides the NAc with a spatial signal that is both environmentally anchored and biased toward behaviourally relevant locations.

\subsection{Speed-inhibited and lick-excited cells are enriched in dHPC$\to$NAc neurons}

Linear regression of calcium events against velocity bins identified speed-modulated neurons in both populations. While speed-excited neurons showed similar proportions across populations, speed-inhibited (deceleration-related) neurons were significantly overrepresented among dHPC$\to$NAc neurons (Figure~4). This enrichment is functionally relevant because mice must decelerate as they approach the hidden reward zone.

Separately, analysis of peri-lick time histograms around appetitive lick onsets revealed that lick-excited neurons were overrepresented in the dHPC$\to$NAc population (Figure~5A--C). These neurons showed robust activity increases aligned to the onset of appetitive licking during reward consumption.

\subsection{Optogenetic activation of dHPC terminals in NAc induces mouth movement and deceleration}

To test functional relevance, separate C57Bl/6 mice ($n = 16$) received CaMKII-ChR2 ($n = 8$) or control EYFP ($n = 8$) virus in dHPC with fibre optic cannulae targeting NAc. Optogenetic stimulation during treadmill running induced two behaviours in ChR2 mice that were absent in controls: increased mouth/jaw movements (quantified via near-infrared camera motion energy at 75~Hz; Figure~6A--C) and deceleration (reduced running velocity; Figure~6). These effects precisely mirror the enriched coding properties found in dHPC$\to$NAc neurons---lick-related activity and speed-inhibition---providing causal evidence that this projection drives consummatory and approach behaviours.

\subsection{Enhanced conjunctive coding in dHPC$\to$NAc neurons}

Venn diagram analysis of functional overlap revealed that dHPC$\to$NAc neurons were more likely to be conjunctively tuned to multiple variables: place + speed-negative, place + lick, and the full triple conjunction of place + speed + lick were all enriched compared to dHPC$^{-}$ neurons (Figure~7A--B).

A generalized linear model (GLM) predicting each neuron's calcium activity from position, velocity, and appetitive licking confirmed this enhanced conjunctive coding (Figure~8). Likelihood ratio tests comparing the full model against reduced models (each predictor removed in turn) demonstrated that dHPC$\to$NAc neurons had significantly larger contributions from all three predictors and stronger interaction terms. The conjunctive representation provides the NAc with a multiplexed signal that simultaneously encodes where the animal is, how fast it is moving, and whether it is consuming reward---information sufficient to identify proximity to the reward zone.

\section{Discussion}

Our findings reveal that the hippocampal projection to the nucleus accumbens carries an enriched, behaviourally specialized representation that goes well beyond a simple spatial signal. dHPC$\to$NAc neurons are not merely place cells routed to a motor output; they carry enhanced spatial information, reward zone overrepresentation, deceleration signals, and consummatory licking responses, combined through conjunctive coding into a multiplexed output \citep{wikenheiser2016hippocampal}.

The overrepresentation of speed-inhibited neurons is particularly noteworthy in the context of goal-directed navigation. Approach to a hidden reward requires deceleration---the animal must slow down to localize the reward zone precisely. By enriching the NAc input with neurons that increase activity during deceleration, the circuit provides a signal that could directly modulate the transition from navigating to consuming \citep{floresco2015nucleus}.

The optogenetic results provide causal support for this model. Stimulating dHPC terminals in NAc produced exactly the two behaviours predicted by the coding analysis: mouth movements (reflecting the lick-excited population) and deceleration (reflecting the speed-inhibited population). This convergence of correlative and causal evidence strengthens the conclusion that the dHPC$\to$NAc projection transmits a goal-approach signal.

\section{Methods}

\subsection{Animals and surgery}
Six Thy1-GCaMP6s mice (GP4.3, Jackson Lab) were used for imaging experiments, and 16 C57Bl/6 mice for optogenetics. Eighteen mice underwent behavioural training only. Surgery was performed under ketamine (0.13~mg/g) and xylazine (0.01~mg/g) anaesthesia. AAVrg-pgk-Cre was injected into NAc (AP $-$1.3, ML $-$1.0, DV $-$5.0 and $-$4.3~mm; $2 \times 500$~nl at 100~nl/min), and AAV5-hSyn-DIO-mCherry into dHPC (200~nl).

\subsection{Behavioural task}
Mice were trained on a 360-cm self-propelled treadmill belt with six textured zones and a 30-cm reward zone. Training lasted five days with liquid rewards dispensed once per lap upon licking in the reward zone.

\subsection{Two-photon imaging}
Dual-colour two-photon imaging was performed through a 3-mm cranial window over dHPC CA1/prosubiculum. GCaMP6s signals were denoised and deconvolved to extract calcium events. mCherry$^{+}$ neurons were classified as dHPC$\to$NAc.

\subsection{Optogenetics}
AAV2-CaMKII-ChR2-EYFP or control rAAV2-CaMKII-EYFP was injected into dHPC with fibre optic cannulae targeting NAc. Orofacial movements were tracked with near-infrared cameras at 75~Hz and quantified as motion energy.

\subsection{Statistical analysis}
Training progression was analyzed with repeated-measures ANOVA (Greenhouse-Geisser corrected when appropriate). Place cell proportions were compared with chi-squared tests. Speed tuning used linear regression on velocity bins. The GLM used likelihood ratio tests for model comparison.

\bibliography{references}

\end{document}
